\SourceLeaf{242}%
\Needspace{6\baselineskip}\begin{center}\textsc{Gauss} to \textsc{Gerling}. Göttingen, 4 February 1844.\end{center}

\mbox{.\ .\ .\ .\ .\ .\ .\ .} That is how it came about that I forgot to answer \textls{another} point in your letter, which I was reminded of just today by a quite accidental circumstance.

The title of the book by my old university friend \textsc{Varkas} (i. e. Wolfgang) von \textsc{Bolyai}, the father, is literally as follows.

Tentamen juventutem studiosam in elementa matheseos purae, elementaris ac sublimioris, methodo intuitiva, evidenti\"aque huic propria, introducendi cum Appendice triplici. Auctore Professore Matheseos et Physices Chemiaeque Publ. Ordinario. Tomus primus. Mar[o]s Vásárhelyini 1832. Typis Collegii Reformatorum per Josephum, et Simeonem Kali de felsö Vist. LXVI and 502 pages 8\textsuperscript{vo}. Tomus secundus 1833. XVI and 400 pages, many copperplates. I do not find the author's name given.

The one appendix attached to the first volume bears the title

Appendix scientiam spatii absolute veram exhibens: a veritate aut falsitate Axiomatis XI Euclidei (a priori haud unquam decidenda) independentem; adjecta ad casum falsitatis, quadratura circuli geometrica. Auctore \textsc{Johanne} \textsc{Bolyai} de eadem\HistoricalNote{*)}{*)}{\emph{of that ilk}, as it is usually put in \textsc{Walter Scott}. — [\emph{See also p.} 225.]}, Geometrarum in exercitu Caesareo Regio Austriaco Castrensium Capitaneo.

This appendix is separately paginated (p. 1---26). Since neither place nor year is given on the title page, I presume that this work, \textls{which is precisely the one in question}, reached the public only as an appendix to the father's work.

In recent decades, a Russian (\textsc{Lobachevsky}, state councillor and professor in Kazan) has taken a similar path. He calls non-Euclidean geometry imaginary geometry (just as your former colleague [\textsc{Schweikart}] called it astral geometry) and has published many very extensive papers about it in \textls{Russian}, mostly in the \gausscyr{Записки Казанскаго Университета}, Memoirs of the University of Kazan, and some in separate pamphlets. I think I possess all of them, but I have put off reading them closely
